\section{Limitations and Future Work}

This article presents a safety-net framing for spoken drone interaction in shared physical spaces. We argue that nearby speech should be processed on the drone side and in real time before it can matter to drone behavior. We develop three design choices for this framing: focusing on speech as the safety-net input, keeping the first speech-processing step local to the drone, and producing timely categories for emergency-related speech, movement-related speech, and unsupported input. We then describe an implementation on a XIAO ESP32-S3 Sense board and summarize evidence from rotor-noisy speech processing, baseline comparison, embedded runtime, and an emergency-related speech scene.

The first limitation is the small category space. The current safety net uses \emph{emergency}, \emph{movement}, and \emph{unknown} because these categories make the main interaction boundary easy to explain: urgent speech, ordinary motion-related speech, and input that should not be acted on are kept separate. Real deployments may need more categories, such as hold, land, return, resume, cancel, request attention, or ask for help. Expanding the vocabulary should be done carefully. Each new category should correspond to a clear interaction need near the drone, not merely to a larger command list.

A related direction is to extend simple speech processing toward stronger intent analysis and recognition. Prior spoken-language-understanding work shows that systems can move from audio toward semantic or intent decisions without relying only on full transcripts~\cite{haghani2018audio_semantics,lugosch19_interspeech,kuo20_interspeech,le2022deliberation}. This direction is useful because different utterances can express the same underlying need. For example, ``stop,'' ``wait,'' and ``do not move closer'' may all indicate an interruption or safety-related request. Similar drone-relevant intent may also appear across languages, accents, and speaker groups, even when the surface words differ. Future systems could therefore learn language- or user-specific expressions while preserving a shared set of drone-relevant intents. This would let the safety net support more diverse speakers without assuming that everyone uses the same words.

Cross-user and cross-language robustness should be studied together with the acoustic setting. Emergency speech may be shouted, repeated, clipped, or masked by rotor noise. A visitor, child, worker, or first responder may use different vocabulary and different levels of urgency. Future datasets should therefore include more speakers, languages, environments, microphone placements, and drone operating conditions. The goal is not only higher recognition accuracy, but a clearer understanding of when the speech-processing module should produce an actionable category and when it should remain conservative.

The implementation can also be improved. On-device speech processing must balance latency, memory, energy, microphone placement, and model size. Future implementations may use better audio front ends, adaptive noise handling, model compression, or hardware-aware scheduling to make the pipeline more reliable across drone platforms. These optimizations matter because the safety-net idea depends on speech being available locally and quickly, especially when a phone or cloud service is unavailable.

Finally, speech should eventually be combined with context and other sensing modalities. Speech alone may be ambiguous when the speaker is far away, not addressing the drone, or masked by noise. Visual cues, gesture, localization, flight state, proximity to people, and task context can help decide whether nearby speech is relevant and how conservatively the drone system should respond.
